% !TEX root = ../aa_SoM_Chapters.tex

%% HARRIS: Chapter 10

% ö = \%c3\%b6
% ä = \%C3\%A4

\xdef\sectiontitle{\IIsectiontitle} %aiheuttama

%%%%%%%%%%%%%%%
\begin{frame}[label=2_1_a]%[label=2_0_TOC1]
\frametitle{\texorpdfstring{\culine[\sectiontitleulinecolor]{\sectiontitle}}{\sectiontitle} }

\vspace{0.2cm}
\begin{itemize}[leftmargin=0.5cm,itemsep=3mm]
    \item[2.1] \hyperlink{2_1_a}{\CDAlert[blue]{\IIaaSubsectiontitle}}
    \item[2.2] \hyperlink{2_2_a}{\IIaSubsectiontitle}
    \item[2.3] \hyperlink{2_3_a}{\IIbSubsectiontitle}	
    \item[2.4] \hyperlink{2_4_a}{\IIcSubsectiontitle}	
    \item[2.5] \hyperlink{2_5_a}{\IIdSubsectiontitle}
    \item[2.6] \hyperlink{2_6_a}{\IIeSubsectiontitle}
    \item[2.7] \hyperlink{2_7_a}{\IIfSubsectiontitle}
    \item[2.8] \hyperlink{2_8_a}{\IIgSubsectiontitle}
    \item[2.9] \hyperlink{2_9_a}{\IIhSubsectiontitle}
    \item[2.10] \hyperlink{2_10_a}{\IIiSubsectiontitle}
\end{itemize}

\endofslide \end{frame}
%%%%%%%%%%%%%%%

\xdef\sectiontitle{\IIaaSubsectiontitle} 

%%%%%%%%%%%%%%%
\begin{frame}%[label=2_1_a]
\frametitle{\texorpdfstring{\culine[\sectiontitleulinecolor]{\sectiontitle}}{\sectiontitle} }

\begin{itemize}[itemsep=4mm]
	\item The \CDAlert[red]{range of distances} associated with atoms\appear{, their nuclei}\appear{, and molecules \CDAlert[red]{is huge!}}
	\item For example, common molecular sizes \appear{are of the order of 0.1–1\;nm:}
\item[] \begin{itemize} \semismall
	\item nitrogen molecule $d=0.37 \mathrm{\;nm}$
	\item water molecule $d=0.3 \mathrm{\;nm}$
\end{itemize}

\item The \CDAlert[blue]{bond lengths} in solids/molecules are $\sim$0.1–0.2\;nm \appear{= 1–2\;å}

\item The \CDAlert[blue]{Bohr radius is $a_{0}=0.0529$\;nm} \appear{(ground state orbit of hydrogen atom)}

\item The size of \CDAlert[fuchsia]{atomic nucleus} is ${\sim}1-10\;\mathrm{fm}=1 ... 10 \times 10^{-6}\;\mathrm{nm}$\appear{, a value dramatically smaller than typical bond lengths}
\implies{ \Alert[tuniblue]{Bonding is mostly associated with outer, i.e.,} \CDAlert[tuniblue]{valence electrons} }
\end{itemize}
%
%\vspace{2mm}
%%
%\begin{tikzpicture}[remember picture,overlay] % anchor=south
%    \node (A) [yshift=1.45*\figYshift,xshift=\figXshift, anchor=east] at (current page.east) {\includegraphics[page=29,scale=\figscale]{Figures_booklet/Geom/Tikz_main_geom_suom_kolumbus.pdf}};
%\end{tikzpicture}
%

\endofslide 
%\endofsection
\end{frame}
%%%%%%%%%%%%%%%

%%%%%%%%%%%%%%%
\begin{frame}
\frametitle{\texorpdfstring{\culine[\sectiontitleulinecolor]{\sectiontitle}}{\sectiontitle} }

\begin{itemize}
	\item For atoms close to each other\appear{, their electrons start seeing a \CDAlert[fuchsia]{'multiatom potential energy'}}
	
	\item Fig.~1 shows four lowest energy states\appear{/wave functions} \appear{in two neighbouring $L$-sized \CDAlert[blue]{finite square wells}} \appear{in situations of atoms with varying separation $a$} 
	
\item When atoms come closer\appear{, both wave functions} \appear{and energy levels are affected.} \appear{Decreasing the distance, the energy levels will spread further apart}
\end{itemize}

\xdef\loc{south east}
\begin{tikzpicture}[remember picture,overlay] % anchor=south
    \node (A) [yshift=-0.25*\figYshift,xshift=\figXshift, anchor=\loc] at (current page.\loc) {\includegraphics[page=1,width=10.5cm]{Figures_booklet/SOM_10_Bonding_figures/SOM_Bonding_Figure_1.png}}; %scale=\figscale. % 8cm max height
\end{tikzpicture}

\endofslide 
%\endofsection
\end{frame}
%%%%%%%%%%%%%%%

%%%%%%%%%%%%%%%
\begin{frame}
\frametitle{\texorpdfstring{\culine[\sectiontitleulinecolor]{\sectiontitle}}{\sectiontitle} }

\begin{itemize}
	\item If there were \CDAlert[red]{three neighbouring finite energy wells}\appear{, then with smaller separation} \appear{the case $n=1$ would have \Alert[red]{three molecular states} } 
\end{itemize}
% 6.5 cm = midway, 10 cm = 3/4 
\begin{minipage}{6.2cm}
\begin{itemize}
	  \item[] with slightly different energy levels\appear{, which at large separation would converge to three isolated atom $n=1$ states}

\item When generalised to a \Alert[blue]{system of $N$ square wells} close to each other%
\appear{\xdef\loc{south east}%
\begin{tikzpicture}[remember picture,overlay]% anchor=south
    \node (A) [yshift=-0.25*\figYshift,xshift=\figXshift, anchor=\loc] at (current page.\loc) {\includegraphics[page=1,height=7cm]{Figures_booklet/SOM_10_Bonding_figures/SOM_Bonding_Figure_2.png}};%scale=\figscale. % 8cm max height
\end{tikzpicture}%
}\appear{, then with decreasing separation distance between the square wells}\appear{, the energy levels will spread and form \CDAlert[blue]{$N$-state energy bands}}
\end{itemize}
\end{minipage}




%\endofslide 
\endofsection
\end{frame}
%%%%%%%%%%%%%%%

